\section{Milestone W0: cubic attachment and quartic normalization}\label{bpW:w0n}

Throughout, Standing assumption~\ref{bpW:w0:mass} ($m>0$) and the conventions of \S\ref{bpW:w0} are in force. The
results of this part are:
(a) the coefficient system \eqref{bpW:w0:coeffsys} is equivalent to the PDE (Proposition~\ref{bpW:w0n:pde-equiv});
(b) $\Gq:\ell^{1,1}\to\ell^{1,2}$ is a bounded cubic polynomial with an explicit constant
(Theorem~\ref{bpW:w0n:cubic-bound});
(c) closed form of the quartic coefficients (Lemma~\ref{bpW:w0n:quartic-coeff});
(d) the action coefficients $A,B$ of the partial quartic Birkhoff normal form for $J=\{1,2\}$
(Theorem~\ref{bpW:w0n:AB});
(e) the nondegeneracy conditions of the source for every $m>0$ (Theorem~\ref{bpW:w0n:nondeg});
(f) the Duffing consistency check (Proposition~\ref{bpW:w0n:duffing}).
All proofs are complete except where a step is explicitly cited to a pinned source and labelled as such; those
citations belong to later milestones (W1, W3) and are listed at the end (\S\ref{bpW:w0n:ledger}).

\subsection{Reconstruction}\label{bpW:w0n:recon}

\begin{lemma}[Reconstruction is bounded into $C^k$]\label{bpW:w0n:recon-bounded}
Let $a>0$, $p\in\R$, $k\in\N$. For $c\in\ell^{a,p}_\C$,
\[\sum_{j\ge1}\abs{c_j}\,j^k\le M_{a,p,k}\norm{c}_{a,p},\qquad M_{a,p,k}^2=\sum_{j\ge1}j^{2k-2p}e^{-2aj}<\infty.\]
Consequently the series $\sum_jc_j\phi_j$ and all its termwise $x$-derivatives of order $\le k$ converge absolutely and
uniformly on $[0,\pi]$; $c\mapsto\sum_jc_j\phi_j$ is a bounded linear map $\ell^{a,p}\to C^k([0,\pi])$ of norm
$\le\sqrt{2/\pi}\,M_{a,p,k}$, its value vanishes at $x=0,\pi$, and $\int_0^\pi(\sum_jc_j\phi_j)\phi_l\,dx=c_l$. In
particular $U_m$, $V_m$ and $q\mapsto U_m(\lambda q)$ are bounded linear maps $\ell^{a,p}\to C^k([0,\pi])$ for every
$p\in\R$, because $\lambda_j^{\pm1/2}\le(1+m)^{1/4}j^{1/2}$ and multiplication by $(\lambda_j^r)$ maps $\ell^{a,p}$
into $\ell^{a,p-\abs r}$ for $\abs r\le 2$, $r\ge0$, with norm $\le(1+m)^{r/2}$.
\end{lemma}

\begin{proof}
Cauchy--Schwarz: $\sum\abs{c_j}j^k=\sum(\abs{c_j}j^pe^{aj})(j^{k-p}e^{-aj})\le\norm c_{a,p}M_{a,p,k}$, and
$M_{a,p,k}<\infty$ because $a>0$. Since $\abs{\partial_x^r\phi_j}\le\sqrt{2/\pi}\,j^r$, the Weierstrass M-test gives
uniform convergence of the derivative series of order $r\le k$, hence termwise differentiation. Boundary values and
the coefficient formula follow termwise from $\phi_j(0)=\phi_j(\pi)=0$ and orthonormality. The bound on
$\lambda_j^r$ is $j^2+m\le(1+m)j^2$.
\end{proof}
\source{P96-W Lemma~1, p.~278, first two sentences (asserted there without the estimate).}

\subsection{The quartic coefficients}\label{bpW:w0n:coeff}

\begin{lemma}[Closed form, (c)]\label{bpW:w0n:quartic-coeff}
For $i,j,k,l\ge1$ let $g_{ijkl}=\int_0^\pi\phi_i\phi_j\phi_k\phi_l\,dx$. Then
\begin{multline*}g_{ijkl}=\frac{1}{2\pi}\Bigl([i+l=j+k]+[i+k=j+l]+[i+j=k+l]\\-[i=j+k+l]-[j=i+k+l]-[k=i+j+l]-[l=i+j+k]\Bigr),\end{multline*}
where $[\cdot]$ is the Iverson bracket. In particular $g$ is symmetric, $g_{ijkl}=0$ unless
$\pm i\pm j\pm k\pm l=0$ for some choice of signs, $g_{iiii}=\frac{3}{2\pi}$, and $g_{iijj}=\frac1\pi$ for $i\ne j$.
Equivalently, for $l\ge1$ and $i,j,k\ge1$,
\begin{equation}\label{bpW:w0n:g-signed}
g_{ijkl}=\frac1{2\pi}\sum_{\varepsilon\in\{\pm1\}^3}s(\varepsilon)\,[\,l=\varepsilon_1i+\varepsilon_2j+\varepsilon_3k\,],
\qquad s(\varepsilon)=(-1)^{1+\#\{t:\varepsilon_t=-1\}},
\end{equation}
so that $\abs{g_{ijkl}}\le\frac1{2\pi}\sum_\varepsilon[\,l=\varepsilon_1i+\varepsilon_2j+\varepsilon_3k\,]$.
\end{lemma}

\begin{proof}
\proofstep{Product to sum}\label{stepW:w0n:prod}
$\sin\alpha\sin\beta=\frac12(\cos(\alpha-\beta)-\cos(\alpha+\beta))$ gives
\[\phi_i\phi_j\phi_k\phi_l=\frac{1}{\pi^2}\bigl(\cos(i-j)x-\cos(i+j)x\bigr)\bigl(\cos(k-l)x-\cos(k+l)x\bigr).\]
\proofstep{Cosine integrals}\label{stepW:w0n:cos}
For integers $A,B$: $\int_0^\pi\cos Ax\cos Bx\,dx=\frac\pi2\bigl([A=B]+[A=-B]\bigr)$ (this includes $A=B=0$, value
$\pi$). Hence
\begin{multline*}g_{ijkl}=\frac1{2\pi}\Bigl([i-j=k-l]+[i-j=l-k]-[i-j=k+l]-[i-j=-k-l]\\-[i+j=k-l]-[i+j=l-k]+[i+j=k+l]+[i+j=-k-l]\Bigr).\end{multline*}
The last bracket vanishes for positive indices; rewriting the others gives the stated form. The special values are
read off: for $(i,i,i,i)$ the three positive brackets hold and no negative one; for $(i,i,j,j)$, $i\neq j$, exactly
$[i+l=j+k]$ and $[i+k=j+l]$ hold (a negative bracket would force one index to be the sum of the other three, e.g.\
$i=i+2j$, impossible).
\proofstep{Signed form}\label{stepW:w0n:signed}
The seven relations, solved for $l$, are $l=-i+j+k$, $l=i-j+k$, $l=i+j-k$ (sign $+$, one minus),
$l=i-j-k$, $l=-i+j-k$, $l=-i-j+k$ (sign $-$, two minuses), $l=i+j+k$ (sign $-$, no minus); the eighth pattern
$l=-i-j-k$ is impossible for $l\ge1$ and may be added with either sign. This is \eqref{bpW:w0n:g-signed}.
\end{proof}
\source{P96-W (10)--(11), p.~281 (only the selection rule and $G_{iijj}$ are stated, ``by an elementary calculation
-- see [6]''); KP96 p.~13 (computation of $G_{ijij}$). The weighted coefficient of P96-W is
$G_{ijkl}=g_{ijkl}/\sqrt{\lambda_i\lambda_j\lambda_k\lambda_l}$; its (11) agrees with the lemma.}

\subsection{The cubic gradient}\label{bpW:w0n:gradient-sec}

\begin{lemma}[$G$ and its gradient as absolutely convergent sums]\label{bpW:w0n:gradient}
Let $a>0$, $p\ge0$, $q\in\ell^{a,p}$ and write $\hat q_j=q_j/\sqrt{\lambda_j}$. Then
\[G(q)=\frac14\sum_{i,j,k,l\ge1}g_{ijkl}\hat q_i\hat q_j\hat q_k\hat q_l,\qquad
\Gq(q)_l=\frac{1}{\sqrt{\lambda_l}}\sum_{i,j,k\ge1}g_{ijkl}\hat q_i\hat q_j\hat q_k,\]
both sums converging absolutely; $G$ is G\^ateaux differentiable along every coordinate with
$\partial G/\partial q_l=\Gq(q)_l$; and $\int_0^\pi U_m(q)^3\phi_l\,dx=\sqrt{\lambda_l}\,\Gq(q)_l$.
\end{lemma}

\begin{proof}
By Lemma~\ref{bpW:w0n:recon-bounded} with $k=0$, $S=\sum_j\abs{\hat q_j}<\infty$ and $U_m(q)=\sum\hat q_j\phi_j$
uniformly. Since $\abs{\phi_j}\le\sqrt{2/\pi}$, the quadruple series $\sum\hat q_i\hat q_j\hat q_k\hat q_l\phi_i\phi_j\phi_k\phi_l$
is dominated by $(2/\pi)^2S^4$ uniformly in $x$; integrate termwise (dominated convergence for series) to obtain both
formulas and their absolute convergence. Differentiating $\frac14\int U_m(q+te_l)^4$ at $t=0$ under the integral
(the integrand is a polynomial in $t$ with coefficients continuous in $x$) gives
$\int U_m(q)^3\phi_l/\sqrt{\lambda_l}$, which is $\Gq(q)_l$ by Definition~\ref{bpW:w0:Gq} and the triple-series
formula.
\end{proof}

\begin{lemma}[Weighted convolution algebra]\label{bpW:w0n:algebra}
Let $a\ge0$ and $s\ge1$. For $c,d\in\ell^{a,s}_b$,
\[\norm{c*d}_{a,s}\le K_s\norm c_{a,s}\norm d_{a,s},\qquad K_s=2^{s}\Bigl(1+2\zeta(2s)\Bigr)^{1/2}.\]
In particular $K_{3/2}^2=8(1+2\zeta(3))$. The convolution sum converges absolutely for each $l$.
\end{lemma}

\begin{proof}
Let $w_l=\jap l^se^{a\abs l}$, $C_l=w_l\abs{c_l}$, $D_l=w_l\abs{d_l}$.
\proofstep{Weight inequality}\label{stepW:w0n:weight}
For $i,l\in\Z$: $e^{a\abs l}\le e^{a\abs i}e^{a\abs{l-i}}$ and $\jap l\le\jap i+\jap{l-i}$ (if $\abs l\le1$ the left
side is $1$; otherwise $\abs l\le\abs i+\abs{l-i}$). By convexity of $t\mapsto t^s$, $(\alpha+\beta)^s\le
2^{s-1}(\alpha^s+\beta^s)$, so $w_l\le2^{s-1}(w_ie^{a\abs{l-i}}+e^{a\abs i}w_{l-i})\le 2^{s-1}(w_iw_{l-i}/\jap{l-i}^s
+w_iw_{l-i}/\jap i^s)$.
\proofstep{Young and Cauchy--Schwarz}\label{stepW:w0n:young}
Hence $w_l\abs{(c*d)_l}\le2^{s-1}\sum_i\bigl(C_i\,\abs{d_{l-i}}e^{a\abs{l-i}}+\abs{c_i}e^{a\abs i}D_{l-i}\bigr)
=2^{s-1}\bigl((C*D')_l+(C'*D)_l\bigr)$ with $C'_i=C_i/\jap i^s$, $D'_i=D_i/\jap i^s$. Young's inequality
$\norm{X*Y}_{\ell^2}\le\norm X_{\ell^2}\norm Y_{\ell^1}$ and $\norm{D'}_{\ell^1}\le\norm D_{\ell^2}
(\sum_i\jap i^{-2s})^{1/2}=\norm D_{\ell^2}(1+2\zeta(2s))^{1/2}$ give the claim; the same bounds give absolute
convergence of each convolution sum.
\end{proof}
\source{P96-W Lemma~2, p.~279 and Appendix~A, pp.~294--295; KP96 Lemma~2 p.~12 and Appendix~A, pp.~31--32 (constant
``depending only on $s$'', not made explicit there). This lemma is the W1 node \texttt{bpW:w1:algebra}; it is proved
here because (b) needs it.}

\begin{theorem}[Cubic gradient, (b)]\label{bpW:w0n:cubic-bound}
Let $a\ge0$, $p\ge\frac12$, and define the symmetric trilinear map $T$ on $(\ell^{a,p}_\C)^3$ by
\[T(q,q',q'')_l=\frac{1}{\sqrt{\lambda_l}}\sum_{i,j,k\ge1}g_{ijkl}\,\hat q_i\hat q'_j\hat q''_k\qquad
(\hat q_j=q_j/\sqrt{\lambda_j}).\]
The sums converge absolutely, $T$ is a bounded trilinear map into $\ell^{a,p+1}_\C$, real on real arguments, and
\[\norm{T(q,q',q'')}_{a,p+1}\le C_*\norm q_{a,p}\norm{q'}_{a,p}\norm{q''}_{a,p},\qquad
C_*=\frac{K_{p+1/2}^2}{\pi}.\]
For $a>0$, $\Gq(q)=T(q,q,q)$ (Lemma~\ref{bpW:w0n:gradient}). Consequently, for $a=p=1$,
$\Gq:\ell^{1,1}\to\ell^{1,2}$ is real analytic (indeed a continuous cubic polynomial), and
\[\norm{\Gq(q)}_{1,2}\le C_m\norm q_{1,1}^3,\qquad C_m=\frac{8(1+2\zeta(3))}{\pi}\approx8.67.\]
The constant is independent of $m>0$ (we only use $\lambda_j\ge j$).
\end{theorem}

\begin{proof}
\proofstep{Even majorants}\label{stepW:w0n:majorant}
Given $q\in\ell^{a,p}_\C$ define the even two-sided sequence $x_i=\abs{q_{\abs i}}/\sqrt{\lambda_{\abs i}}$ for
$i\neq0$, $x_0=0$; similarly $x',x''$. Using $\lambda_i\ge i$,
\[\norm x_{a,p+1/2}^2=2\sum_{i\ge1}\frac{i}{\lambda_i}\,i^{2p}e^{2ai}\abs{q_i}^2\le2\norm q_{a,p}^2.\]
\proofstep{Selection rule bound}\label{stepW:w0n:selection}
By \eqref{bpW:w0n:g-signed}, for $l\ge1$,
\[\sum_{i,j,k\ge1}\abs{g_{ijkl}}\,\abs{\hat q_i\hat q'_j\hat q''_k}\le\frac1{2\pi}\sum_{i,j,k\ge1}\sum_{\varepsilon}
[\,l=\varepsilon_1i+\varepsilon_2j+\varepsilon_3k\,]\,x_ix'_jx''_k\le\frac{1}{2\pi}(x*x'*x'')_l,\]
since $(i,j,k,\varepsilon)\mapsto(\varepsilon_1i,\varepsilon_2j,\varepsilon_3k)$ is injective into $(\Z\setminus0)^3$
and $x$ is even. Thus the sum defining $T$ converges absolutely.
\proofstep{Smoothing}\label{stepW:w0n:smoothing}
With $X=x*x'*x''$ (even, by evenness of the factors) and $l/\lambda_l\le1$,
\[\norm{T}_{a,p+1}^2\le\sum_{l\ge1}\frac{l^{2p+2}e^{2al}}{\lambda_l}\frac{X_l^2}{4\pi^2}
\le\frac1{4\pi^2}\sum_{l\ge1}l^{2p+1}e^{2al}X_l^2\le\frac{1}{8\pi^2}\norm X_{a,p+1/2}^2.\]
By Lemma~\ref{bpW:w0n:algebra} with $s=p+\frac12\ge1$, $\norm X_{a,p+1/2}\le K_s^2\norm x\norm{x'}\norm{x''}\le
K_s^22^{3/2}\norm q\norm{q'}\norm{q''}$. Hence $\norm T_{a,p+1}\le(2\sqrt2\pi)^{-1}2^{3/2}K_s^2\prod\norm\cdot=
(K_s^2/\pi)\prod\norm\cdot$.
\proofstep{Analyticity}\label{stepW:w0n:analytic}
Trilinearity and symmetry are clear from the formula and the symmetry of $g$. A bounded symmetric trilinear map
$T$ between complex Banach spaces defines the continuous homogeneous cubic polynomial $q\mapsto T(q,q,q)$, which
is analytic on the whole space, with the finite power series $T(q_0+h)^{3}$ expanded by multilinearity (in Lean:
a \Lean{ContinuousMultilinearMap} is analytic). It maps real sequences to real sequences. For $a>0$ the identification
with $\Gq$ is Lemma~\ref{bpW:w0n:gradient}.
\end{proof}
\source{P96-W Lemma~3, pp.~279--280 (with $O(\cdot)$ instead of a constant) and (13)--(14), p.~284; KP96 p.~16.
The source deduces analyticity from coordinatewise analyticity and local boundedness [P96-W ref.~9, Appendix~A];
the polynomial structure makes that principle unnecessary here.}
\sourcegap{P96-W p.~280 asserts $\norm{\mathcal Sq}_{a,s+1/2}\le\norm q_{a,s}$ in the two-sided norm. The correct
bound is $\le\sqrt2\norm q_{a,s}$ (Step~\ref{stepW:w0n:majorant}); the printed one fails for $q=e_1$, $a=s=0$,
$m=1$. Constants only; recorded as new item N-1 in the errata check.}

\subsection{PDE \texorpdfstring{$\leftrightarrow$}{<->} coefficient system}\label{bpW:w0n:pde}

\begin{proposition}[Equivalence, (a)]\label{bpW:w0n:pde-equiv}
Let $I\subset\R$ be an open interval and $q,p:I\to\ell^{1,1}$.
\begin{enumerate}[label=(\roman*)]
\item (Coefficients $\Rightarrow$ PDE.) Suppose $q,p$ are continuous into $\ell^{1,1}$, continuously differentiable
  into $\ell^{1,0}$, and satisfy \eqref{bpW:w0:coeffsys} as an identity in $\ell^{1,0}$. Then
  $u(t,x)=U_m(q(t))(x)$ is a classical solution of \eqref{bpW:w0:eq} on $I$ (Definition~\ref{bpW:w0:nlw}) with
  $u_t(t,\cdot)=V_m(p(t))$.
\item (PDE $\Rightarrow$ coefficients.) Suppose $u(t,x)=U_m(q(t))(x)$ is a classical solution of
  \eqref{bpW:w0:eq} on $I$ and $u_t(t,\cdot)=V_m(p(t))$ for every $t$. Then every $q_j,p_j$ is $C^1$ on $I$ and
  \eqref{bpW:w0:coeffsys} holds coordinatewise.
\end{enumerate}
In both directions the cubic term is realized exactly: for $q\in\ell^{1,1}$,
\begin{equation}\label{bpW:w0n:force}
U_m(q)^3=\sum_{l\ge1}\sqrt{\lambda_l}\,\Gq(q)_l\,\phi_l\quad\text{uniformly on }[0,\pi],\qquad
-\partial_x^2U_m(q)+mU_m(q)=U_m(\lambda^2q),
\end{equation}
where $(\lambda^2q)_j=\lambda_j^2q_j$ and the right sides converge in $C^0$ (first) and $C^0$ (second, as
$\lambda^2q\in\ell^{1,-1}$).
\end{proposition}

\begin{proof}
\proofstep{Force identity}\label{stepW:w0n:force}
Let $f=U_m(q)^3\in C([0,\pi])$. By Lemma~\ref{bpW:w0n:gradient}, $\int f\phi_l=\sqrt{\lambda_l}\Gq(q)_l$; by
Theorem~\ref{bpW:w0n:cubic-bound} and Lemma~\ref{bpW:w0n:recon-bounded} (the sequence
$(\sqrt{\lambda_l}\Gq(q)_l)$ lies in $\ell^{1,1}$), $S=\sum_l\sqrt{\lambda_l}\Gq(q)_l\phi_l$ converges uniformly to a
continuous $S$ with the same sine coefficients as $f$. Since $(\phi_l)_{l\ge1}$ is a complete orthonormal system of
$L^2(0,\pi)$, $f=S$ almost everywhere, hence everywhere by continuity. The second identity is termwise
differentiation (Lemma~\ref{bpW:w0n:recon-bounded} with $k=2$) and $-\phi_j''+m\phi_j=\lambda_j^2\phi_j$.
\proofstep{Time regularity}\label{stepW:w0n:time}
In (i), $\dot q=\lambda p$ is continuous into $\ell^{1,0}$, and since $p$ is $C^1$ into $\ell^{1,0}$ and
multiplication by $\lambda$ is bounded $\ell^{1,0}\to\ell^{1,-1}$, $\dot q$ is $C^1$ into $\ell^{1,-1}$ with
$\ddot q=\lambda\dot p=-\lambda^2q-\lambda\Gq(q)$. As $U_m:\ell^{1,-1}\to C^2([0,\pi])$ is bounded linear,
$t\mapsto u(t,\cdot)$ is $C^2$ into $C^2([0,\pi])$; therefore $u,u_t,u_{tt},u_x,u_{xx}$ exist and are jointly
continuous, $u_t=U_m(\lambda p)=V_m(p)$ and $u_{tt}=U_m(\ddot q)$.
\proofstep{Identification}\label{stepW:w0n:ident}
By \eqref{bpW:w0n:force}, $U_m(\lambda^2q)=-u_{xx}+mu$ and $U_m(\lambda\Gq(q))=\sum_l\sqrt{\lambda_l}\Gq(q)_l\phi_l=u^3$;
hence $u_{tt}=u_{xx}-mu-u^3$. The boundary conditions hold by Lemma~\ref{bpW:w0n:recon-bounded}.
\proofstep{Converse}\label{stepW:w0n:converse}
In (ii), $\int_0^\pi u(t,x)\phi_l(x)\,dx=q_l(t)/\sqrt{\lambda_l}$ and $\int_0^\pi u_t\phi_l=\sqrt{\lambda_l}p_l$.
Differentiation under the integral sign ($u_t,u_{tt}$ continuous on compact sets) gives
$\dot q_l=\sqrt{\lambda_l}\int u_t\phi_l=\lambda_lp_l$ and $\sqrt{\lambda_l}\dot p_l=\int u_{tt}\phi_l$. Two
integrations by parts, using $u(t,0)=u(t,\pi)=0=\phi_l(0)=\phi_l(\pi)$, give $\int u_{xx}\phi_l=-l^2\int u\phi_l$, so
$\int u_{tt}\phi_l=-\lambda_l^2q_l/\sqrt{\lambda_l}-\sqrt{\lambda_l}\Gq(q)_l$ by Lemma~\ref{bpW:w0n:gradient}. Divide by
$\sqrt{\lambda_l}$.
\end{proof}
\source{P96-W Lemma~1, p.~278, and (7), p.~278.}
\sourcegap{P96-W Lemma~1 compresses termwise differentiation, completeness of the sine system and the identification
of $f(u)$ with its sine series into one paragraph and states only (i) for real-analytic curves. The proof above
supplies the steps; completeness of $(\phi_l)$ in $L^2(0,\pi)$ is a standard external fact (Fourier basis of
$L^2$ on the circle applied to odd extensions). The all-time real-analytic version used by the headline is W2.}

\subsection{Quartic nonresonance}\label{bpW:w0n:nonres-sec}

\begin{lemma}[Qualitative nonresonance of all quartic terms]\label{bpW:w0n:nonres}
Let $m>0$, let $a,b,c,d\ge1$ with $g_{abcd}\neq0$, and let $\sigma\in\{\pm1\}^4$. If
$\sigma_1\lambda_a+\sigma_2\lambda_b+\sigma_3\lambda_c+\sigma_4\lambda_d=0$, then the indices with $\sigma=+1$ and
those with $\sigma=-1$ form equal multisets (in particular there are two of each).
\end{lemma}

\begin{proof}
By Lemma~\ref{bpW:w0n:quartic-coeff}, $g_{abcd}\ne0$ implies a signed relation $\pm a\pm b\pm c\pm d=0$; in
particular each index is at most the sum of the other three.
\proofstep{Four or three equal signs}\label{stepW:w0n:31}
All signs equal is impossible ($\lambda>0$). If three signs agree, say $\lambda_a+\lambda_b+\lambda_c=\lambda_d$,
then $d\le a+b+c$, so $\lambda_d\le\lambda_{a+b+c}$, while $\lambda_x+\lambda_y>\lambda_{x+y}$ for $x,y\ge1$
(square both sides: $2\sqrt{(x^2+m)(y^2+m)}>2xy-m$); hence $\lambda_a+\lambda_b+\lambda_c>\lambda_{a+b+c}\ge\lambda_d$.
\proofstep{Two against two}\label{stepW:w0n:22}
Let $\lambda_a+\lambda_b=\lambda_c+\lambda_d$ with $\{a,b\}\ne\{c,d\}$ as multisets, and let $f(s)=\sqrt{s^2+m}$.
Up to renaming within and between the two sides, the signed relation is one of:
(1) $a+b=c+d$: for fixed sum, $f(x)+f(y)$ is strictly increasing in $\abs{x-y}$ by strict convexity of $f$, so
equality forces $\abs{a-b}=\abs{c-d}$ and $\{a,b\}=\{c,d\}$;
(2) $a-b=d-c=:t$ (that is, $a+c=b+d$): then $f(b+t)+f(b)=f(c)+f(c+t)$ and $y\mapsto f(y)+f(y+t)$ is strictly
increasing on $\{y\ge1,\ y+t\ge1\}$, so $b=c$, $a=d$;
(3) one index is the sum of the other three, and after renaming it is $a=b+c+d$ (the side containing the largest
index is called the left side). Then $\lambda_a-\lambda_c=\int_c^{b+c+d}f'$ and $\lambda_d-\lambda_b=\int_b^df'$
(nonpositive if $d\le b$). Since $f'>0$ is increasing, $\int_c^{c+b+d}f'>\int_{c+2b}^{c+b+d}f'\ge\int_b^df'$ when
$d>b$ (an interval of the same length $d-b$ shifted right by $c+b>0$). So
$\lambda_a+\lambda_b>\lambda_c+\lambda_d$.
Every signed relation among four positive integers with two signs on each side is of type (1) or (2) after renaming,
and every relation with one sign different is of type (3); this exhausts the cases.
\end{proof}
\begin{remark}\label{bpW:w0n:nonres-remark}
The lemma is qualitative and covers \emph{all} indices. It does not give a uniform lower bound: divisors such as
$\lambda_{j+1}-\lambda_j-\lambda_{k+1}+\lambda_k$ tend to $0$, which is why P96-W (p.~276) says no complete quartic
normal form is available. The quantitative bound for quartic terms with at least one index in $\{1,\dots,n\}$,
$\abs{\lambda'_i+\lambda'_j+\lambda'_k+\lambda'_l}\ge cm/\sqrt{n^2+m}^{\,3}$, is P96-W Lemma~4, p.~283 (proof
pp.~285--286); it belongs to W3 and is cited, not reproved, here.
\end{remark}

\subsection{Partial quartic Birkhoff normal form: the action coefficients}\label{bpW:w0n:birkhoff}

Write $\hat\lambda=(\lambda_j)$ and, in the complex coordinates of Definition~\ref{bpW:w0:z},
$q_j=(z_j+\bar z_j)/\sqrt2$. Then
\begin{equation}\label{bpW:w0n:Gz}
G=\frac14\sum_{i,j,k,l}\frac{g_{ijkl}}{\sqrt{\lambda_i\lambda_j\lambda_k\lambda_l}}q_iq_jq_kq_l
=\frac1{16}\sum_{i,j,k,l}\frac{g_{ijkl}}{\sqrt{\lambda_i\lambda_j\lambda_k\lambda_l}}
(z_i+\bar z_i)(z_j+\bar z_j)(z_k+\bar z_k)(z_l+\bar z_l).
\end{equation}
Expanding gives a sum of monomials $z^\alpha\bar z^\beta$ with $\abs\alpha+\abs\beta=4$. Along the linear flow
($\dot z_j=-i\lambda_jz_j$) such a monomial evolves by $e^{-i\avg{\alpha-\beta,\hat\lambda}t}$; it is
\emph{resonant} if $\avg{\alpha-\beta,\hat\lambda}=0$. Let $J=\{1,2\}$, $N=\{j\ge3\}$, and split
$G=G^{\rm tan}+G^{\rm nor}$, where $G^{\rm nor}$ collects the monomials all of whose indices lie in $N$.

\begin{definition}[Partial quartic normal form]\label{bpW:w0n:Qdef}
$\bar G$ is the sum of the resonant monomials of $G^{\rm tan}$ (with their coefficients), and the partial quartic
normal form is $\Lambda+\bar G+G^{\rm nor}$.
\end{definition}

\begin{theorem}[Action coefficients, (d)]\label{bpW:w0n:AB}
For $J=\{1,2\}$ and every $m>0$, with $I=(I_1,I_2)$, $Z=(I_j)_{j\ge3}$ and $I_j=\abs{z_j}^2$,
\[\bar G=\frac12\avg{AI,I}+\avg{BI,Z},\qquad
A_{ii}=\frac{9}{8\pi\lambda_i^2},\quad A_{12}=A_{21}=\frac{3}{2\pi\lambda_1\lambda_2},\quad
B_{ji}=\frac{3}{2\pi\lambda_j\lambda_i}\ (j\ge3,\ i\in J).\]
Equivalently $A=\frac{3}{8\pi}D(4\mathbf 1\mathbf 1^{\mathsf T}-\mathrm{Id})D$ with $D=\operatorname{diag}(\lambda_1^{-1},\lambda_2^{-1})$,
and $(BI)_j=\frac{3}{2\pi\lambda_j}\bigl(\frac{I_1}{\lambda_1}+\frac{I_2}{\lambda_2}\bigr)$. More generally, the full
angle average of $G$ is
\begin{equation}\label{bpW:w0n:avgG}
\avg G=\frac{3}{16\pi}\Bigl(4\Bigl(\sum_{j\ge1}\frac{I_j}{\lambda_j}\Bigr)^2-\sum_{j\ge1}\frac{I_j^2}{\lambda_j^2}\Bigr).
\end{equation}
\end{theorem}

\begin{proof}
\proofstep{Resonant equals averaged}\label{stepW:w0n:resavg}
By Lemma~\ref{bpW:w0n:nonres}, a monomial $z^\alpha\bar z^\beta$ of \eqref{bpW:w0n:Gz} with nonzero coefficient is
resonant iff $\alpha=\beta$, i.e.\ it is $\abs{z_a}^2\abs{z_b}^2=I_aI_b$. With $z_j=\sqrt{I_j}e^{-i\varphi_j}$ the
average over independent angles of $z^\alpha\bar z^\beta$ is $\delta_{\alpha\beta}I^\alpha$. Hence $\bar G$ is the
part of $\avg G$ (average over all angles, computed on finitely supported $q$ and extended by the absolutely
convergent formula) consisting of the terms that contain at least one tangential action.
\proofstep{Averages of monomials}\label{stepW:w0n:moments}
With $q_j=\sqrt{2I_j}\cos\varphi_j$: $\avg{q_iq_jq_kq_l}=0$ unless the indices pair up; $\avg{q_i^4}=4I_i^2\cdot\frac38
=\frac32I_i^2$ and $\avg{q_i^2q_j^2}=I_iI_j$ for $i\neq j$. Among ordered quadruples, $(i,i,i,i)$ occurs once and the
multiset $\{i,i,j,j\}$ ($i<j$) six times.
\proofstep{Summation}\label{stepW:w0n:sum}
Therefore, with $y_j=I_j/\lambda_j$ and Lemma~\ref{bpW:w0n:quartic-coeff},
\[\avg G=\frac14\Bigl(\sum_i\frac{3}{2\pi}\cdot\frac32y_i^2+\sum_{i<j}6\cdot\frac1\pi\,y_iy_j\Bigr)
=\frac{9}{16\pi}\sum_iy_i^2+\frac{3}{2\pi}\sum_{i<j}y_iy_j,\]
which is \eqref{bpW:w0n:avgG}. Matching with $\frac12\sum_{i,j}A_{ij}I_iI_j=\frac12\sum_iA_{ii}I_i^2+\sum_{i<j}A_{ij}I_iI_j$
on $J$ and with $\sum_{j\ge3,i\in J}B_{ji}I_iZ_j$ on the mixed terms gives the coefficients.
\end{proof}

\begin{proposition}[Formal partial Birkhoff transformation]\label{bpW:w0n:formal-birkhoff}
Let $F$ be the quartic polynomial whose coefficient on each nonresonant monomial $z^\alpha\bar z^\beta$ of
$G^{\rm tan}$ with coefficient $c_{\alpha\beta}$ is $c_{\alpha\beta}/(-i\avg{\alpha-\beta,\hat\lambda})$, and zero
otherwise. Then $F$ is real ($\overline{F}=F$ under $z\leftrightarrow\bar z$), $X_\Lambda F=G^{\rm tan}-\bar G$ where
$X_\Lambda F=\frac{d}{dt}F$ along the linear flow, and formally
\[H\circ\Phi^1_F=\Lambda+\bar G+G^{\rm nor}+(\text{terms of degree}\ \ge6),\]
where $\Phi^1_F$ is the time-one map of $X_F$. The quartic part of the normal form is independent of the choice of
$F$ among solutions of $X_\Lambda F=G^{\rm tan}-\bar G$.
\end{proposition}

\begin{proof}
The divisors are nonzero by Lemma~\ref{bpW:w0n:nonres}; $X_\Lambda(z^\alpha\bar z^\beta)=-i\avg{\alpha-\beta,\hat\lambda}
z^\alpha\bar z^\beta$ gives the homological identity. Reality: the coefficient of $\bar z^\alpha z^\beta$ is
$\bar c_{\alpha\beta}$ (since $G$ is real) and the divisor changes sign and is conjugated. The Lie series
$H\circ\Phi^1_F=H+\{H,F\}+\frac12\{\{H,F\},F\}+\dots$ with $\{H,F\}=X_FH=-X_HF$ has quartic part
$G-X_\Lambda F=G^{\rm nor}+\bar G$. Any two solutions $F$ differ by a resonant quartic $F_0$, and $X_\Lambda F_0=0$.
\end{proof}
\source{P96-W Main Proposition, pp.~281--282; proof pp.~282--284, (12) and the Taylor formula on p.~283. The
analytic realization of $\Phi^1_F$ on a neighbourhood of $0$ in $E_{a,p}$, the smoothing of $X_F$, and the bound
$O(\norm z^6)$ on the remainder are W3 (they use P96-W Lemma~4 and (13)--(14)). The \emph{coefficients} printed
there, $\bar G_{ij}=\frac6\pi\frac{4-\delta_{ij}}{\lambda_i\lambda_j}$, are $16$ times those of
Theorem~\ref{bpW:w0n:AB}; see Remark~\ref{bpW:w0n:compare}.}

\begin{remark}[Comparison with the printed coefficients]\label{bpW:w0n:compare}
P96-W p.~282 expands $G$ as $\frac14\sum G_{ijkl}(z_i+\bar z_i)\cdots(z_l+\bar z_l)$; by \eqref{bpW:w0n:Gz} the
prefactor is $\frac1{16}$. The next line rewrites this as $\sum' G_{ijkl}w_iw_jw_kw_l$ \emph{without} any prefactor,
and the multiplicity count of p.~284 ($\bar G_{ij}=24G_{iijj}$, $\bar G_{ii}=12G_{iiii}$) is correct for that line.
Two factors of $4$ are therefore lost; with the prefactor $\frac1{16}$ the same count gives exactly $A$ and $B$ above.
The factor is common to $A$ and $B$, so the nondegeneracy arguments of P96-W Lemmas~5--6 are unaffected
(Remark~\ref{bpW:w0:harmless}).
\end{remark}

\subsection{Nondegeneracy for every \texorpdfstring{$m>0$}{m>0}}\label{bpW:w0n:nondeg-sec}

The Cantor Manifold Theorem used by P96-W (p.~288, assumption~A) requires, with $\omega(I)=\lambda_J+AI$ and
$\Omega(I)=\lambda_N+BI$,
(A$_1$) $\det A\ne0$; (A$_2$) $\avg{l,\lambda_N}\ne0$; (A$_3$) $\avg{k,\omega(I)}+\avg{l,\Omega(I)}\not\equiv0$ as
an affine function of $I$, for all $(k,l)\in\Z^2\times\Z^{N}$ with $1\le\abs l\le2$ ($\abs l=\sum\abs{l_j}$).

\begin{theorem}[Nondegeneracy, (e)]\label{bpW:w0n:nondeg}
For $J=\{1,2\}$ and every $m>0$:
\begin{enumerate}[label=(\roman*)]
\item $\det A=-\dfrac{63}{64\pi^2\lambda_1^2\lambda_2^2}=-\dfrac{63}{64\pi^2(1+m)(4+m)}\neq0$;
\item $\avg{l,\lambda_N}\neq0$ for $1\le\abs l\le2$;
\item (A$_3$) holds.
\end{enumerate}
\end{theorem}

\begin{proof}
\proofstep{Determinant}\label{stepW:w0n:det}
$\det A=(\tfrac3{8\pi})^2\lambda_1^{-2}\lambda_2^{-2}\det\begin{psmallmatrix}3&4\\4&3\end{psmallmatrix}=
-\tfrac{63}{64\pi^2}\lambda_1^{-2}\lambda_2^{-2}$. (For general $n$: $\det(4\mathbf1\mathbf1^{\mathsf T}-\mathrm{Id})=
(-1)^{n-1}(4n-1)$.)
\proofstep{Pure normal}\label{stepW:w0n:A2}
$\avg{l,\lambda_N}$ is $\pm\lambda_a$, $\pm2\lambda_a$, $\pm(\lambda_a+\lambda_b)$ or $\pm(\lambda_a-\lambda_b)$ with
$a\neq b$ in $N$; all are nonzero since $\lambda>0$ is strictly increasing.
\proofstep{Reduction of (A$_3$)}\label{stepW:w0n:A3red}
Let $b=(\lambda_1^{-1},\lambda_2^{-1})$ and $s=\sum_{j\ge3}l_j/\lambda_j$. Then
$A=\frac3{8\pi}(4bb^{\mathsf T}-D^2)$ and $B^{\mathsf T}l=\frac{3}{2\pi}s\,b$, so the linear part of the affine map
is $\frac{3}{8\pi}\bigl(4(\avg{b,k}+s)b-D^2k\bigr)$. It vanishes iff $k_i/\lambda_i=c$ for $i=1,2$ with
$c=4(\avg{b,k}+s)$; then $\avg{b,k}=2c$ and $s=-\frac74c$. The constant part is
$c(\lambda_1^2+\lambda_2^2)+\avg{l,\lambda_N}$.
\proofstep{The case $\abs l=1$}\label{stepW:w0n:l1}
$l=\sigma e_\nu$, $\nu\ge3$: $s=\sigma/\lambda_\nu$, $c=-4\sigma/(7\lambda_\nu)$, and the constant part vanishes iff
$7\lambda_\nu^2=4(\lambda_1^2+\lambda_2^2)$. Then $k_i^2=c^2\lambda_i^2=\frac{4\lambda_i^2}{7(\lambda_1^2+\lambda_2^2)}
\in(0,\frac47)$, not the square of an integer. So (A$_3$) holds.
\proofstep{The case $\abs l=2$}\label{stepW:w0n:l2}
$\abs s\le2/\lambda_3<\frac23$, so $\abs c<\frac{8}{21}$, and $k_2-k_1=c(\lambda_2-\lambda_1)$ with
$0<\lambda_2-\lambda_1<1$; integrality forces $k_1=k_2$, hence $c=0$, hence $s=0$. But $s\ne0$ for $\abs l=2$: $s=0$
would need $l=e_a-e_b$ with $\lambda_a=\lambda_b$, $a\neq b$. So the linear part does not vanish.
\end{proof}
\source{P96-W assumption~A p.~288; Lemma~5 pp.~291--292 ((A$_1$), with the printed matrix; the printed
$(1-4n)\prod\lambda_j^{-2}$ agrees after scaling); Lemma~6 pp.~292--293 ((A$_3$) for admissible $J$). Our proof of
(iii) is the specialization of Lemma~6 to $J=\{1,2\}$; the second branch of the source's $\abs l=2$ argument
($\abs{\avg{w,l}}\ge(4n-1)/(4n-4)$) cannot occur for $J=\{1,2\}$ since $\abs{\avg{w,l}}\le2/\lambda_3<1$.}

\begin{remark}[Spectral asymptotics with the corrected $B$]\label{bpW:w0n:kappa}
$\Omega_j(I)=\lambda_j+\frac{3}{2\pi\lambda_j}\avg{b,I}=j+\frac{m_I}{j}+O(j^{-3})$ with
$m_I=\frac m2+\frac{3}{2\pi}\avg{b,I}$, so $(\Omega_i-\Omega_j)/(i-j)=1+O(j^{-2})$ uniformly for bounded $I$: the
exponent $\kappa=2$ of P96-W p.~291 is unchanged by the correction (only $v=\frac{24}\pi b$ becomes $\frac{3}{2\pi}b$).
\end{remark}

\subsection{Duffing consistency check}\label{bpW:w0n:duffing-sec}

\begin{proposition}[One-mode frequency shift, (f)]\label{bpW:w0n:duffing}
Restrict \eqref{bpW:w0:eq} to $u=x(t)\phi_j(x)$ with only mode $j$ excited and drop the forcing of other modes
(which is $O(x^3)$ and does not enter the first-order frequency shift). The coefficient equation is the Duffing
equation $\ddot x+\lambda_j^2x+\varepsilon x^3=0$ with $\varepsilon=g_{jjjj}=\frac{3}{2\pi}$, whose
Lindstedt--Poincar\'e frequency at amplitude $a$ is $\lambda_j+\frac{3\varepsilon a^2}{8\lambda_j}+O(a^4)$. With
$q_j=\sqrt{\lambda_j}x$ and $I_j=\frac12\lambda_ja^2$ this is
\[\lambda_j+\frac{9}{16\pi\lambda_j}a^2=\lambda_j+\frac{9}{8\pi\lambda_j^2}I_j=\lambda_j+A_{jj}I_j,\]
the frequency $\partial_{I_j}(\lambda_jI_j+\frac12A_{jj}I_j^2)$ of the normal form. The printed P96-W value
$\bar G_{jj}=\frac{18}{\pi\lambda_j^2}$ would predict a shift $16$ times larger.
\end{proposition}

\begin{proof}
$\int(x\phi_j)^3\phi_j=x^3g_{jjjj}$ gives the Duffing coefficient. The Lindstedt shift $3\varepsilon a^2/(8\omega_0)$ is
classical (e.g.\ by averaging $\varepsilon x^3$ against $\cos$ over a period: $\frac{2}{T}\int\varepsilon
a^3\cos^4=\frac34\varepsilon a^3$). Along a solution $x=a\cos\varphi$, $q_j=\sqrt{\lambda_j}a\cos\varphi$ and
$p_j=\dot q_j/\lambda_j$, so to leading order $I_j=\frac12(q_j^2+p_j^2)=\frac12\lambda_ja^2$. Substitute.
\end{proof}
\begin{remark}
The relation between $I_j$ and the amplitude holds to leading order only; the agreement checked is of the
coefficient of the first-order shift, which is what the normal form determines. \texttt{scripts/wave/w0\_checks.py}
also integrates the Duffing equation numerically for $m\in\{0.1,1,10\}$ and confirms
$\abs{\omega_{\rm num}-(\lambda_1+A_{11}I_1)}=O(I_1^2)$.
\end{remark}

\subsection{Ledger: proved, cited, source gaps}\label{bpW:w0n:ledger}

\begin{itemize}
\item \emph{Proved here}: Lemmas~\ref{bpW:w0n:recon-bounded}, \ref{bpW:w0n:quartic-coeff},
  \ref{bpW:w0n:gradient}, \ref{bpW:w0n:algebra}, \ref{bpW:w0n:nonres}; Theorems~\ref{bpW:w0n:cubic-bound},
  \ref{bpW:w0n:AB}, \ref{bpW:w0n:nondeg}; Propositions~\ref{bpW:w0n:pde-equiv}, \ref{bpW:w0n:formal-birkhoff}
  (formal level), \ref{bpW:w0n:duffing}.
\item \emph{External standard facts}: completeness of $(\phi_l)$ in $L^2(0,\pi)$; analyticity of continuous
  multilinear maps; Young's inequality; $\zeta(3)<\infty$; the classical Lindstedt shift.
\item \emph{Cited, for later milestones}: quantitative divisor bound (P96-W Lemma~4, W3); analytic realization of
  the Birkhoff map and its remainder (P96-W Main Proposition, W3).
\item \emph{Source gaps / normalization defects}: P96-W p.~280 constant (N-1); P96-W Lemma~1 compression; P96-W
  p.~282/284 factor $16$ (E-PDE-W-1); P96-W pp.~287, 290 action/recovery/sign conventions (E-PDE-W-2).
\end{itemize}
